\Needspace{10\baselineskip}
\subsection{A.157. A finite resolution by fibre joins over an arbitrary real closed field}

\begin{saresult}{A.157}{Proved under stated hypotheses}{A universal finite-join resolution and a relative metric consequence}
Let \(M\) be any o-minimal expansion of a real closed field. All sets, maps, homotopies, compactness assertions, and products below are definable over \(M\). Here compact means definably compact. Let \(f:K\to Y\) be a continuous surjection between compact definable sets. If \(Y\ne\varnothing\) and \(d=\dim Y\), there are a compact definable set \(R\), a continuous surjection \(q:R\to Y\), an integer \(1\le m\le d+1\), and a continuous map
\[
 \ell:R\longrightarrow J_m(f),\qquad \pi\ell=q,
 \tag{157.1}
\]
where \(J_m(f)\) is the labelled \(m\)-fold fibrewise join of \(f\). The map \(q\) is universally a definable homotopy equivalence: for every continuous map \(T\to Y\) from a compact definable set, \(R\times_YT\to T\) is a definable homotopy equivalence.

If \(Z\subset Y\) is closed and \(b:Z\to J_a(f)\) is a continuous section of its join projection, the resolution can be enlarged to a compact \(R'\) containing \(Z\) as a closed subspace, with \(q'|_Z=\operatorname{id}_Z\), a universal homotopy equivalence \(q':R'\to Y\), and a map
\[
 \ell':R'\longrightarrow J_{m+a}(f),\qquad
 \pi\ell'=q',\qquad \ell'|_Z=b
 \tag{157.2}
\]
(the last equality uses the last \(a\) slots). In particular, for \(Y=P=[0,1]^k\), \(Z=\partial P\), there is a continuous definable \(w:P\to J_{m+a}(f)\) with \(w|_{\partial P}=b\), and \(\pi w\) is definably homotopic to \(\operatorname{id}_P\) relative to \(\partial P\).

For the following metric consequence impose an additional hypothesis, which is not supplied by the topological assertion. Suppose \(C_p\), \(p\in P\), is a bounded Hausdorff-continuous definable family of nonempty compact subsets of a fixed \(M^n\). Suppose there is a compact definable \(K\) and a Lipschitz surjection \(\phi:K\to P\) such that \(s\mapsto C_{\phi(s)}\) is Hausdorff-Lipschitz. If the boundary family \(C|_{\partial P}\) is Hausdorff-Lipschitz, there is a definable Lipschitz map \(g:P\to P\), equal to the identity on \(\partial P\), such that \(p\mapsto C_{g(p)}\) is Hausdorff-Lipschitz. The convex homotopy from \(\operatorname{id}_P\) to \(g\) fixes \(\partial P\). This consequence holds over every \(M\) as above, conditional on the displayed Lipschitz-resolution hypothesis; it does not silently extend the real-field theorem of Kocel--Cynk, Pawlucki, and Valette.
\end{saresult}

\begin{proof}
We use only finite constructions. The required all-field topology is provided by Barton--Commelin, \emph{Model categories for o-minimal geometry}, arXiv:2108.11952: Proposition 6.17, printed p.~29, gives existence and pullback stability of closed-by-proper pushouts; Propositions 8.5--8.6, pp.~47--48, give the homotopy extension property for closed embeddings of compact definable sets and the resulting cofibration category, whose weak equivalences are definable homotopy equivalences; Proposition 9.28, p.~60, gives a finite closed splitting filtration for every surjective definable map; Proposition 9.31, p.~63, identifies the proper pushout square used below. No ordinary finite subcover or subdivision by an infinitesimal mesh is invoked.

\textbf{Finite mapping-cylinder facts.}
For maps \(A\xrightarrow{a}E\) and \(A\xrightarrow{b}F\) between compact definable sets, write
\[
 D(E\leftarrow A\to F)=
 E\ \cup_a\ (A\times[0,1])\ \cup_b\ F
 \tag{157.3}
\]
for their double mapping cylinder. It is a compact definable space and can be realized as a closed bounded definable set. For completeness, after embedding the three sets in finite coordinate spaces, map the cylinder point \((x,t)\) to
\[
 ((1-t)a(x),\,t b(x),\,t(1-t)x,\,t).
 \tag{157.4}
\]
Include the endpoint copies \((e,0,0,0)\), \(e\in E\), and \((0,z,0,1)\), \(z\in F\). For \(0<t<1\), the third coordinate recovers \(x\); at the two ends exactly the identifications in (157.3) occur. The image is compact and definable, and the induced proper quotient is the required realization. Equivalently, (157.3) is formed by two closed-by-proper pushouts.

We use two finite homotopy-pushout facts. An objectwise homotopy equivalence between two spans induces a homotopy equivalence of their double mapping cylinders. Also, if \(A\to E\) is a closed embedding, the canonical map from (157.3) to the ordinary pushout \(E\cup_AF\) is a homotopy equivalence. These are the gluing consequences of the cofibration category in Proposition 8.6: replace a leg by its mapping cylinder, whose free endpoint is a closed cofibration, and apply the gluing lemma for cofibrations. Thus these facts concern finite definable mapping cylinders, not ordinary topological realizations. They also apply after any compact definable base change. Indeed, each of the two endpoint pushouts in (157.3) is pullback-stable by Proposition 6.17, and a closed embedding remains a closed embedding after base change.

\textbf{The labelled join.}
Embed \(K\) in some \(M^N\). Define \(J_m(f)\) as the image of
\[
 \{(s_1,\ldots,s_m)\in K^m:f(s_1)=\cdots=f(s_m)\}
          \times\Delta^{m-1}
 \longrightarrow M^{m(N+1)},\qquad
 (s,t)\longmapsto(t_i,t_i s_i)_{i=1}^m.
 \tag{157.5}
\]
It is compact and definable. A zero-weight slot forgets its \(s_i\). Conversely, any such inactive slots can be completed using one active slot, so this model is exactly the labelled fibrewise join. Its projection \(\pi\) takes the common value \(f(s_i)\) of an active slot. The quotient construction makes \(\pi\) continuous. Adding zero-weight slots gives the usual inclusions among the joins.

\textbf{Induction constructing a universal resolution.}
If \(Y=\varnothing\), then \(K=\varnothing\); use \(R=\varnothing\) and the unique maps. If \(Y\) is nonempty of dimension zero, it is finite; choose one point of each fibre and take \(R=Y\), \(q=\operatorname{id}\), \(m=1\). This case is universal under every base change.

For the induction step, a definable choice section of \(f\) is continuous on a relatively open set \(U\subset Y\) whose closed complement \(C\) has dimension less than \(\dim Y\). This is precisely the dimension-dropping step in Proposition 9.28. Set \(K_C=f^{-1}(C)\) and
\[
 E=\overline{s(U)}^{\,K}\cup K_C,\qquad
 A=E\cap K_C=K_C.
 \tag{157.6}
\]
The addition of \(K_C\) avoids needing a separate density assertion. The restriction \(E\to Y\) is proper and onto, and is a homeomorphism over \(U\): the graph of the continuous section is closed over the relatively open set \(U\). Thus the square
\[
 \begin{matrix}
 A&\longrightarrow&E\\
 \downarrow&&\downarrow\\
 C&\longrightarrow&Y
 \end{matrix}
 \tag{157.7}
\]
is a pullback and a pushout by Proposition 9.31. The upper arrow is a closed embedding.

If \(C=\varnothing\), simply take \(R=E\), \(q=f|_E\), \(m=1\). Otherwise apply the induction hypothesis to \(K_C\to C\), obtaining
\[
 q_C:R_C\to C,\qquad
 \ell_C:R_C\to J_{m_C}(K_C\to C),\qquad
 m_C\le\dim C+1.
\]
Put \(A'=A\times_CR_C\), and form
\[
 R=D(E\leftarrow A'\to R_C).
 \tag{157.8}
\]
All three pieces map to \(Y\), with a common value on each cylinder line; hence they define \(q:R\to Y\). The projection \(A'\to A\) is a homotopy equivalence by the \emph{universal} induction hypothesis. Therefore the map of spans from (157.8) to \(E\leftarrow A\to C\) is an objectwise homotopy equivalence. The first mapping-cylinder fact compares their double mapping cylinders, and the second compares the latter with (157.7). It follows that \(q\) is a homotopy equivalence.

Here is the full universality check. Given compact \(T\to Y\), append a subscript \(T\) to each base change. Pullback stability of the two endpoint pushouts identifies \(R\times_YT\), as a definable space, with
\[
 D(E_T\leftarrow A'_T\to (R_C)_T).
\]
The maps \(A'_T\to A_T\) and \((R_C)_T\to C_T\) are base changes of \(q_C\), so are homotopy equivalences. The map \(A_T\to E_T\) is still a closed embedding, and (157.7) remains a pushout with lower-right corner \(T\). The same two finite mapping-cylinder facts therefore show \(R\times_YT\to T\) is a homotopy equivalence. No density or section assertion about \(U_T\) is required.

To define \(\ell\), use the point \(e\in E\subset K\) in the first slot on the \(E\)-piece, and use \(\ell_C(r)\) in the remaining \(m_C\) slots on the \(R_C\)-piece. On a cylinder point \((a,r,t)\), with \(f(a)=q_C(r)\), give the first slot weight \(1-t\), and multiply every weight of \(\ell_C(r)\) by \(t\). These points lie in one common fibre. The formula agrees with the endpoint values, so descends to a continuous definable map \(R\to J_{m_C+1}(f)\) over \(Y\). Since \(m_C+1\le d+1\), this completes the induction.

\textbf{A prescribed section on a closed boundary.}
Given \(Z\subset Y\) and \(b\) as in (157.2), set \(R_Z=R\times_YZ\). It is closed in \(R\), and \(R_Z\to Z\) is a homotopy equivalence by universality. Form
\[
 R'=D(R\leftarrow R_Z\to Z).
 \tag{157.9}
\]
The inclusion \(R\to R'\) is a homotopy equivalence: compare (157.9) to the double mapping cylinder \(D(R\leftarrow R_Z\to R_Z)\), which retracts onto \(R\), using the homotopy equivalence \(R_Z\to Z\). Thus \(q':R'\to Y\) is a homotopy equivalence by two-out-of-three. The identical argument after compact base change shows it is universal. The endpoint copy \(Z\subset R'\) is a closed embedding and \(q'|_Z=\operatorname{id}\).

Map \(R\) to the first \(m\) join slots by \(\ell\), and \(Z\) to the last \(a\) slots by \(b\). On \(R_Z\times I\), multiply the first weights by \(1-t\) and the last weights by \(t\). Again all active slots have the same base value, so this defines \(\ell'\) over \(Y\).

When \(Y=P\) and \(Z=\partial P\), choose a definable homotopy inverse \(v:P\to R'\), together with a homotopy \(v q'\simeq\operatorname{id}_{R'}\). Its restriction to the endpoint copy of \(\partial P\) is a homotopy from \(v|_{\partial P}\) to that inclusion. Apply the definable homotopy extension property of \(\partial P\subset P\) to change \(v\) to \(u:P\to R'\) with \(u|_{\partial P}\) exactly the endpoint inclusion. Set \(w=\ell'u\). Then \(w|_{\partial P}=b\), while \(\pi w=q'u\) fixes \(\partial P\). Because \(P\) is convex, the straight homotopy from \(\operatorname{id}_P\) to \(q'u\) remains in \(P\) and fixes \(\partial P\). We do not assert that \(w\) is a strict interior lift of \(\operatorname{id}_P\).

\textbf{The conditional metric assertion.}
Use maximum norms. Add to \(K\) a positively separated tagged copy of \(\partial P\); map that copy to \(\partial P\) by the identity and assign its original family \(C_p\). Boundedness, positive separation, and the assumed boundary Lipschitz estimate make both the resulting projection \(\phi^+:K^+\to P\) and its pulled-back compact-set family globally Lipschitz. The boundary section into one join slot is affine in \(p\) in the coordinates (157.5).

For any finite \(m\), both the projection and the pulled-back family on \(J_m(\phi^+)\) are Lipschitz in its actual Euclidean coordinates. Here is the estimate. Let \(M_0=\max_{s\in K^+}\|s\|_\infty\). For two join points \(z=(t_i,v_i)\), \(z'=(t_i',v_i')\), put \(\delta=\|z-z'\|_\infty\). If \(\delta<1/(2m)\), choose \(t_i\ge1/m\). Then \(t_i'>0\), and
\[
 \left\|\frac{v_i}{t_i}-\frac{v_i'}{t_i'}\right\|_\infty
 \le m(1+M_0)\delta.
 \tag{157.10}
\]
Evaluate the projection or the compact-set family on this active slot. For \(\delta\ge1/(2m)\), use the bounded diameter of \(P\), respectively a bound for all Hausdorff distances of the family. This proves a global Lipschitz estimate without selecting slots continuously.

The preceding relative construction supplies a continuous definable \(w:P\to J_m(\phi^+)\) with the literal affine boundary value \(b\), for some finite \(m\le k+2\). We now make \(w\) Lipschitz without any subdivision argument. Apply the total-variation coordinate-clock construction of A.154 to the finite-dimensional vector map \((\operatorname{id}_P,w)\). That construction applies to any finite target dimension: its proof uses only a finite maximum formula for variation, uniform o-minimal monotonicity, and the estimate
\[
 |\operatorname{Var}(h)-\operatorname{Var}(h')|
       \le 2N\|h-h'\|_\infty
\]
for scalar line functions with at most \(N\) monotonicity intervals. It gives a definable homeomorphism \(\chi:P\to P\) preserving every coordinate face such that \(w\chi\) is globally Lipschitz.

In this application \(\chi\) fixes the boundary pointwise. At each stage, a processed coordinate line contained in a boundary face has \(w=b\), which is affine. Its summed coordinate variation is therefore \(ct\) for a constant \(c\) on that line, and the normalized clock \((t+ct)/(1+c)\) is exactly \(t\). A boundary face which fixes the processed coordinate is fixed because the interval clock fixes its two endpoints. Induction preserves the affine boundary values at every stage. This argument is valid over every real closed field.

Finally put \(g=\pi w\chi\). The join estimates make \(g\) Lipschitz and make \(C_g\) Hausdorff-Lipschitz. Its boundary value is the identity. The convex relative homotopy completes the proof. If the parameter cube has dimension zero the assertion is immediate, and empty spatial families, when separately allowed, have their unique trivial regularization.
\end{proof}

\paragraph{Source and remaining scope.}
The finite topological resolution is independent of Kocel--Cynk--Pawlucki--Valette and is proved over arbitrary o-minimal real closed fields. Their inspected original manuscript of 5 September 2012 explicitly fixes the ordered field of real numbers in its introduction; therefore the existence of the initial Lipschitz resolution in the metric assertion is retained as a hypothesis outside the already checked real-field or real-model cases. The argument neither assumes a triangulation of the original map nor invokes ordinary Smale theory over a non-Archimedean field.
